% concatenated from main.tex
%%%%%%%%%%%%%%%%%%%%%%%%%%%%%%%%%%%%%%%%%%%%%%%%%%%%%%%%%%%%%%%%%%%%%%%%%%%%%%%%
%2345678901234567890123456789012345678901234567890123456789012345678901234567890
%        1         2         3         4         5         6         7         8

\documentclass[letterpaper, 10 pt, conference]{ieeeconf}  % Comment this line out
                                                          % if you need a4paper
%\documentclass[a4paper, 10pt, conference]{ieeeconf}      % Use this line for a4
                                                          % paper

\IEEEoverridecommandlockouts                              % This command is only
                                                          % needed if you want to
                                                          % use the \thanks command
\overrideIEEEmargins
% See the \addtolength command later in the file to balance the column lengths
% on the last page of the document



% The following packages can be found on http:\\www.ctan.org
%\usepackage{graphics} % for pdf, bitmapped graphics files
%\usepackage{epsfig} % for postscript graphics files
%\usepackage{mathptmx} % assumes new font selection scheme installed
%\usepackage{times} % assumes new font selection scheme installed
%\usepackage{amsmath} % assumes amsmath package installed
%\usepackage{amssymb}  % assumes amsmath package installed
\usepackage{amsmath, amssymb}
\usepackage{cite}
\usepackage{algorithm}
\usepackage{algpseudocode}
\usepackage{graphicx}
\usepackage{bm}
\newtheorem{theorem}{Theorem}
\newtheorem{proposition}{Proposition}
\newtheorem{corollary}{corollary}

\title{\LARGE \bf
Moral Hazard in LTI Dynamics: A Hypothesis Testing Approach}

%\author{ \parbox{3 in}{\centering Huibert Kwakernaak*
%         \thanks{*Use the $\backslash$thanks command to put information here}\\
%         Faculty of Electrical Engineering, Mathematics and Computer Science\\
%         University of Twente\\
%         7500 AE Enschede, The Netherlands\\
%         {\tt\small h.kwakernaak@autsubmit.com}}
%         \hspace*{ 0.5 in}
%         \parbox{3 in}{ \centering Pradeep Misra**
%         \thanks{**The footnote marks may be inserted manually}\\
%        Department of Electrical Engineering \\
%         Wright State University\\
%         Dayton, OH 45435, USA\\
%         {\tt\small pmisra@cs.wright.edu}}
%}

\author{Jaewon Jeong, Pan-Yang Su, S. Shankar Sastry, and Anil Aswani% <-this % stops a space
\thanks{This work was partially supported by Collaborative Research: Transferable, Hierarchical, Expressive, Optimal, Robust, Interpretable NETworks (THEORINET) under Simons Foundation Award MPS-MODL-00814647 and National Science Foundation (NSF) Grant DMS-2031899, and by the NSF under Grant DGE-2125913.}% <-this % stops a space
\thanks{J. Jeong and A. Aswani are with Industrial Engineering and Operations Research,
        University of California, Berkeley, CA 94720
        {\tt\small \{jaewon\_jeong,aaswani\}@berkeley.edu}}%
\thanks{P.Y. Su and S.S. Sastry are with the Electrical Engineering and Computer Sciences, University of California, Berkeley, CA 94720
        {\tt\small \{pan\_yang\_su,shankar\_sastry\}@berkeley.edu}}%
}


\begin{document}



\maketitle
\thispagestyle{empty}
\pagestyle{empty}


%%%%%%%%%%%%%%%%%%%%%%%%%%%%%%%%%%%%%%%%%%%%%%%%%%%%%%%%%%%%%%%%%%%%%%%%%%%%%%%%
\begin{abstract}
Many incentive design problems must contend with information asymmetries due to non-observation of efficiency (adverse selection) or non-observation of effort (moral hazard). And although a growing body of literature considers incentive design in control systems, the problem of designing incentives for control systems under information asymmetries has been less well-studied. This paper considers a model of moral hazard within control systems. In our model, the control system is described by an (affine) linear time-invariant (LTI) system with process noise. There is an agent who gets to choose (from between two choices) a linear state-feedback controller to apply to the LTI system, with one of the state-feedback controllers having a higher quadratic cost on the control inputs than the other. Our goal is to design a payment scheme that incentivizes the agent to choose the state-feedback controller that minimizes a quadratic cost on system states plus the time-discounted payment amount, subject to the understanding that the agent bears the control cost while being risk-averse with respect to their time-discounted payment. We formulate the problem as a constrained optimization, and prove that for a  payment given after a fixed (but optimizable) time horizon the optimal payment scheme chooses the payment amount using a likelihood ratio hypothesis test. We numerically demonstrate our results by applying the derived optimal payment scheme to two examples: load frequency control (LFC) in power systems and wellness interventions for body weight loss.
\end{abstract}

\section{INTRODUCTION}
%\section{Sy, stem Model}

%Our framework is most closely related to the "Principal-Agent Hypothesis Testing" paradigm introduced by Bates et al. \cite{Bates2022}, which utilizes statistical evidence thresholds to determine agent payouts. However, while Bates et al. focus on an adverse selection model where the agent has private information about a fixed quality parameter, our work addresses the moral hazard problem where the agent’s hidden action—the choice of control law $K_i$—dynamically alters the system's state trajectory. 

Economic aspects of control systems are increasingly important given the growth of applications at the intersection of economics and control, like power \cite{Tan2010,5718175} or transportation \cite{su2024incentive,11298540} systems. As such, an active body of literature within controls considers various problems of incentive design, in which a \emph{principal} would like to design a scheme for payments (or some similar motivating force) to an \emph{agent} to encourage the agent to take actions that are favorable for the principal. Though incentive design in economics focuses on asymmetries of information between participants \cite{Laffont}, the design of incentives in control systems can be challenging even in the absence of information asymmetries. In this work, we propose a principal-agent model for stochastic LTI systems with continuous outputs subject to moral hazard and show that a likelihood ratio test on the continuous output trajectories constitutes an optimal contract for incentivizing the agent's hidden controller choice.
%Our work show that for stochastic LTI systems under moral hazard, the optimal contract to incentivize an agent's hidden controller choice is a likelihood ratio test on the continuous output trajectories.
\subsection{Incentive Design in Control Systems}

Much of the literature on incentive design for control systems focuses on the setting of perfect information, an early example of which is \cite{HO1982167}. In these models, the principal is assumed to have perfect knowledge about the agent and their control actions. This line of work has extended to settings such as bilevel model predictive control (BiMPC) \cite{mintz2018control,mintz2023behavioral} to design incentives in situations similar to what occurs in hierarchical model predictive control (MPC), and this work has recently been generalized to design incentives in more complex hierarchical MPC structures \cite{thirugnanam2025dynamic}.

\subsection{Information Asymmetry in Control Systems}

A closely-related line of literature \cite{ratliff2019perspective} considers incentive design for control systems under various types of information asymmetry. Given the surge of data being generated by modern systems, more recent work has combined ideas from machine learning with control to design incentives under information asymmetry. An example of such are models combining a principal-agent framework with multiarmed bandits with hidden (to the principal) rewards \cite{dogan2023repeated,dogan2023estimating}.

The work most closely related to our model is \cite{Yang2014}, which designs dynamic contracts with real-time and end-time compensation for stochastic control systems by solving an associated Hamilton Jacobi-Bellman (HJB) equation. Our work differs in a few notable ways: The most significant difference is that we propose a simpler contract that consists of a fixed payment time, two payment amounts, and regions of the observation-space corresponding to each payment amount (all of which are optimized over) for a setup where the agent can choose between only two state-feedback controllers. Other differences are that \cite{Yang2014} considers continuous-time nonlinear models, whereas in our setup we focus on discrete-time linear models with linear state-feedback controllers.

\subsection{Contribution and Outline}
% We make two main contributions: 
% \begin{itemize}
%     \item For a model of moral hazard with a stochastic linear time-invariant (LTI) system in discrete-time and two payment levels, we show that the optimal contract depends on a likelihood ratio test. While similar results are known when outputs are discrete \cite{Laffont,saig2023delegated}, to the best of our knowledge, our result is the first to consider continuous outputs, as is the case for LTI systems.
%     \item Under the model, we derive the distribution of the log likelihood ratio statistic under stochastic LTI dynamics, demonstrating that they take the form of generalized sums of chi-squared distributions.
% \end{itemize}
 


%We make two main contributions: For a model of moral hazard with a stochastic linear time-invariant (LTI) system in discrete-time and two payment levels, we show that the optimal contract depends on a likelihood ratio test. Similar results are known when outputs are discrete \cite{Laffont,saig2023delegated}; however, to the best of our knowledge our result is the first to consider continuous outputs, as is the case for LTI systems. We also derive the distribution of the log likelihood ratio statistic under stochastic LTI dynamics, demonstrating that they take the form of generalized sums of chi-squared distributions.

We make two main contributions: For a model of moral hazard with a stochastic linear time-invariant (LTI) system in discrete-time and two payment levels, we show that the optimal contract depends on a likelihood ratio test. Similar results are known when outputs are discrete \cite{Laffont,saig2023delegated}; however, our result  considers continuous outputs. We also derive the distribution of the log likelihood ratio statistic under stochastic LTI dynamics, showing they take the form of generalized sums of chi-squared distributions.

The remainder of this paper is organized as follows. Section \ref{sec:cdp} presents the system model and the contract design problem. In Section \ref{sec:soc}, we theoretically characterize the structure of the optimal contract, and Section \ref{sec:ctoc} provides a practical numerical algorithm to compute the optimal contract. Lastly, Section \ref{sec:ae} presents applications of the proposed contract design to load frequency control in power systems and a wellness program to incentivize body weight loss.


\section{CONTRACT DESIGN PROBLEM}
\label{sec:cdp}
In the problem we consider, there is a \emph{principal} that wishes to control a dynamical system; however, the principal must pay an \emph{agent} to choose control actions on their behalf. The principal incurs a state cost, and the agent incurs a control cost. The agent picks a controller from a choice of two state-feedback policies, with one of the state-feedback policies having a higher expected control cost than the other. The principal's goal is to design a contract that maps any state trajectory to a payment amount, where the payment serves as a cost to the principal and a reward to the agent. %The principal's goal is to design a contract, which is a mapping from any state trajectory to a payment amount, with the payment as a cost to the principal and a reward to the agent. 

\subsection{Control System and Costs}
\subsubsection{System Dynamics}
The system the principal would like to control is described by a discrete-time linear time-invariant (LTI) system with process and measurement noise
\begin{align}
    x_{k+1}&=Ax_{k}+Bu_{k}+w_{k}\\
    y_k &= Cx_k + \epsilon_k
\end{align}
where $x_{k} \in \mathbb{R}^{n}$ is the system state, $y_k \in \mathbb{R}^p$ is the noisy output,  $u_k \in \mathbb{R}^q$ is the control input, the matrices have dimensions $A\in\mathbb{R}^{n\times n}$,  $B\in\mathbb{R}^{n\times q}$, $C\in\mathbb{R}^{p\times n}$, the $w_{k} \sim \mathcal{N}(\mu_{w}, \Sigma_{w})$ with $\mu_w \in \mathbb{R}^n$ and $\Sigma_w \in S^n_+$ (with $S^n_+$ denoting the set of $\mathbb{R}^{n\times n}$ positive semidefinite matrices) is independent and identically distributed (i.i.d.) jointly Gaussian process noise, and $\epsilon_{k} \sim \mathcal{N}(0, \Sigma_{e})$ with $\epsilon_k\in\mathbb{R}^p$ and $\Sigma_e \in S^p_+$ is i.i.d. jointly Gaussian measurement noise. Our model allows the process noise to have non-zero mean $\mu_w$ because this allows us to include affine LTI systems in our results. The initial state $x_0$ has a jointly Gaussian distribution $x_0 \sim \mathcal{N}(\mu_0, \Sigma_{0})$ with $\mu_0 \in \mathbb{R}^n$ and $\Sigma_0 \in S^n_+$. We further assume that $x_0$, $w_t$, $\epsilon_t$ for all $t$ are independent.

In our model, we do \emph{not} require $(A,C)$ to be observable because such an assumption is stronger than needed for our problem. We will instead make a related but weaker assumption about the distributions of the $y_0,y_1,\ldots$. This assumption will appear in the statements of our theorems. 

\subsubsection{Agent Costs}

The agent chooses between two linear state-feedback controllers, one $u_k = K_0 x_k$ which is interpreted as a \emph{low effort} controller, and the other $u_k = K_1 x_k$ which is interpreted as a \emph{high effort} controller. For the matrices $K_0, K_1 \in \mathbb{R}^{q\times p}$, assume  $(A+BK_0)$ and $(A+BK_1)$ are Schur stable. The interpretation of controller effort levels is driven by a time-discounted quadratic control cost
\begin{equation}    
    J^A = \mathbb{E}\Big[ \textstyle\lim_{N \to \infty} \sum_{t=0}^{N}\gamma_a^t \cdot(u_t^\top R u_t + r^\top u_t)\Big],
\end{equation}
where $R \in \mathbb{R}^{q\times q}$ and $r \in \mathbb{R}^q$ are constants of the quadratic cost, and $\gamma_a \in (0,1)$ is a discount factor. Let $J^A_0$ be the value of this cost under the low effort controller $K_0$, and let $J^A_1$ be the value of this cost under the high effort controller $K_1$. We assume the values of the various constant matrices and vectors  in our model are such that $J^A_1 > J^A_0$ to ensure the high versus low effort interpretation of $K_1$ versus $K_0$.

The agent also receives a payment $\pi$ at time $T$ from the principal as an incentive, and the agent is assumed to be risk-averse with respect to this payment amount. We model the time-discounted utility of the payment as $\gamma_a^T\cdot U(\pi)$, where $U: \mathbb{R}\rightarrow\mathbb{R}$ is a mapping from a payment amount to a utility. We assume that $U$ satisfies the following properties, which are standard in risk-aversion models \cite{Laffont}: $U(0) = 0$, $U' > 0$, and $U'' < 0$. 
%with the properties that: it passes through the origin $U(0) = 0$, its first derivative is positive $U' > 0$, and its second derivative is negative $U'' < 0$. This is a standard model for risk-aversion \cite{Laffont}. 
Hence, the agent's total cost is given by $J^A - \gamma_a^T \cdot U(\pi)$. Note that the minus sign implies that positive utility reduces the agent's cost.%(Note the minus sign, which indicates that positive utility reduces agent costs.)

\subsubsection{Principal Costs}

The principal incurs a time-discounted state cost
\begin{equation}
    J^P = \mathbb{E}\Big[ \textstyle\lim_{N \to \infty} \sum_{t=0}^{N}\gamma_p^t \cdot(x_t^\top Q x_t + q^\top x_t)\Big],
\end{equation}
where $Q \in \mathbb{R}^{p\times p}$ and $q \in \mathbb{R}^p$ are constants of the cost, and $\gamma_p \in (0,1)$ is a discount factor. Let $J^P_0$ be the value of this cost under the low effort controller $K_0$, and let $J^P_1$ be the value of this cost under the high effort controller $K_1$. 

The principal makes a payment $\pi$ at time $T$ to the agent as an incentive, and the principal incurs a cost of $\gamma_p^T \cdot \pi$ for this payment. Hence, the principal's total cost is $J^P + \gamma_p^T \cdot \pi$. Note that $\pi$ is a function, and the choice of $\pi$ and $T$ will be formulated as an optimization problem in Section \ref{subsec: contract design problem}.

%(The choice of $\pi$, which is a function, and $T$ will be optimized in the design problem.) Hence, the principal's total cost is given by $J^P + \gamma_p^T \cdot \pi$.


\subsection{Contract Design Problem}
\label{subsec: contract design problem}
In our setup, the principal must make two choices. The first choice is to decide whether or not they would like the agent to exert high effort via a payment. This is a choice because, depending on the coefficient values, the cost of payment that would be required to induce the agent to exert high effort may be higher than the cost reduction the principal gets when the agent exerts high effort. The second choice is that if the principal decides it is beneficial to induce high effort from the agent, then the principal must choose a contract to offer the agent. However, the problem of designing the optimal contract in moral hazard settings typically proceeds backwards \cite{Laffont}: The principal first computes an optimal contract $\pi^*$ given an affirmative decision to induce high effort, and then the principal checks whether it is beneficial to execute this contract to induce high effort. 

In our setting, moral hazard arises because the principal does not directly observe the effort level (i.e., the choice of $K_0$ or $K_1$) of the agent. Instead, we assume the principal takes measurements $y_0, y_1, \ldots$. The mapping between the agent's choice of effort level and the outputs is probabilistic because of the process and measurement noise, and so it is impossible to perfectly identify the agent's choice of effort level from the observations. What this means is that the contract design problem must necessarily be probabilistic.

Before we present the optimal contract design problem, we first formally specify the structure of the class of contracts we consider. Let $\mathbf{Y}_T = (y_1,\ldots,y_T)$. Then a contract is a function $\pi: \mathbf{Y}_T \mapsto \pi(\mathbf{Y}_T) \in \mathbb{R}$ that maps the observed outputs $\mathbf{Y}_T = (y_1,\ldots,y_T)$ over a finite horizon $T$ to a payment $\pi(\mathbf{Y}_T)$, which may be positive or negative. 
%Then a contract is a function $\pi(\mathbf{Y}_T) \in \mathbb{R}$ that inputs a finite horizon $T$ and the observed outputs over this horizon $y_1,\ldots,y_T$ and returns a payment amount, which may be positive or negative. 
For the purpose of practical implementation, it is important for the contract to be easily interpretable by the agent. Consequently, here we further restrict the contract to have the following easily-interpretable structure:
\begin{equation}
\pi(\mathbf{Y}_T) = \begin{cases} \pi_0, & \text{if } \mathbf{Y}_T \notin \mathcal{R} \\ \pi_1, & \text{if } \mathbf{Y}_T \in \mathcal{R} 
\end{cases}
\end{equation}
where the values $T \in\mathbb{Z}_{+}$ and $\pi_0,\pi_1 \in \mathbb{R}$ and the set $\mathcal{R}$ are to be optimized in the contract design problem. The interpretation is that $\mathcal{R}$ is a region for which the state trajectory is more likely to belong to under high effort, whereas the complement is a region for which the state trajectory is more likely to belong to under low effort. The contract then pays $\pi_0$ for low effort, and $\pi_1$ for high effort.

It is now straightforward to formulate the contract design problem assuming the principal wishes to induce high effort. Let $H_1$ be the distributions of $y_0,y_1,\ldots$ that result from our model when the agent chooses high effort $K_1$, and let $H_0$ be the distributions of $y_0,y_1,\ldots$ that result from our model when the agent chooses low effort $K_0$. We assume that $H_0 \neq H_1$. For notational convenience, we define $\alpha(\mathcal{R}) = \mathbb{P}[\mathbf{Y}_T \in \mathcal{R}\ |\ H_0]$ and $\beta(\mathcal{R}) = \mathbb{P}[\mathbf{Y}_T \in \mathcal{R}\ |\ H_1]$. We will drop the argument and write $\alpha,\beta$ when the meaning is clear from the context. The contract design problem is 
\begin{align}
%     \min_T \min_{\pi_0,\pi_1}\min_{\mathcal{L}}\ & J_1^P + \gamma_p^T \cdot \big(\pi_0 \cdot (1-\beta) + \pi_1 \cdot \beta\big) \\
     \min\ & J_1^P + \gamma_p^T \cdot \big(\pi_0 \cdot (1-\beta) + \pi_1 \cdot \beta\big) \label{eqn:cost}\\
     & J_1^A - \gamma_a^T\cdot\big(U(\pi_0) \cdot (1-\beta) + U(\pi_1)\cdot\beta) \leq \nonumber \\
     & \qquad J_0^A - \gamma_a^T\cdot\big(U(\pi_0) \cdot (1-\alpha) + U(\pi_1)\cdot\alpha) \label{eqn:icc}\\
     & J_1^A - \gamma_a^T\cdot\big(U(\pi_0) \cdot (1-\beta) + U(\pi_1)\cdot\beta) \leq J_0^A \label{eqn:pc}
\end{align}
where the optimization is over the horizon $T \in \{1,\ldots,\overline{T}\}$ with a constraint on the maximum  horizon $\overline{T} \in \mathbb{Z}_+$, the payment amounts $\pi_0,\pi_1\in\mathbb{R}$ with $\pi_1 \geq \pi_0$, and the decision region $\mathcal{R} \subseteq \mathbb{R}^{T \cdot p}$. The principal minimizes its expected total cost (\ref{eqn:cost}), subject to a constraint (\ref{eqn:icc}) that ensures the contract is such that the agent's expected cost when exerting high effort is lower than the agent's expected cost when exerting low effort (i.e., an \emph{incentive compatibility} constraint that ensures high effort is incentivized) and a constraint (\ref{eqn:pc}) that ensures the agent incurs a lower expected cost from the contract when exerting higher effort than if they exert low effort and receive no payment (i.e., a \emph{participation} constraint that ensures it is rational for the agent to accept the contract). 

\section{STRUCTURE OF OPTIMAL CONTRACT}
\label{sec:soc}

In this section, we theoretically characterize the mathematical structure of the optimal contract.

\subsection{Characterization of Optimal Contract}
The key result is a theorem that relates our contract design problem to the Neyman-Pearson testing framework. %It is notationally convenient to define an indicator function 
%\begin{equation}
%    \phi(\mathbf{X}_T) = \begin{cases} 1, & \text{if } \mathbf{X}_T \in \mathcal{R}\\
%    0, &\text{if } \mathbf{X}_T \notin\mathcal{R}
%    \end{cases}
%\end{equation}
%for the state trajectory $\mathbf{X}_T$ lying in the decision region $\mathcal{R}$.

\begin{proposition}
\label{prop:optimality}
Suppose $T,\pi_0,\pi_1$ are fixed with $\pi_1 >  \pi_0$, and assume there exists $\mathcal{R}$ such that
\begin{align}
    \alpha &\leq - U(\pi_0)/(U(\pi_1)-U(\pi_0)) \label{eqn:alph}\\
    \beta &= \Delta - U(\pi_0)/(U(\pi_1)-U(\pi_0)) \label{eqn:beth}
\end{align}
for the constant
\begin{equation}
    \Delta := \gamma_a^{-T}\cdot\frac{J^A_1 - J^A_0}{U(\pi_1) - U(\pi_0)}.\label{eqn:Delta}
\end{equation} Then there exists an optimal decision region $\mathcal{R}^*$ that solves (\ref{eqn:cost})--(\ref{eqn:pc}), which takes the form of
\begin{equation}
\label{eqn:lrt}
    \mathcal{R}^* = \big\{\mathbf{Y}_T : L_1(\mathbf{Y}_T) \geq \kappa \cdot L_0(\mathbf{Y}_T)\big\}.
\end{equation}
where $L_0(\cdot)$ and $L_1(\cdot)$ are the (probabilistic) likelihoods of $\mathbf{Y}_T$ under $H_0$ and $H_1$, respectively, and $\kappa \geq 0$ is a scalar.
\end{proposition}

\begin{proof}
Observe that the only decision variable here is $\mathcal{R}$. We can rewrite the principal's cost (\ref{eqn:cost}) as
\begin{equation}
    J_1^P + \gamma_p^T \cdot \pi_0 + \gamma_p^T\cdot\big(\pi_1 -\pi_0\big)\cdot\beta. \label{eqn:cost2}
\end{equation}
Because $\pi_1 > \pi_0$ and $\gamma_p \in (0,1)$, minimizing (\ref{eqn:cost}) subject to (\ref{eqn:icc}), (\ref{eqn:pc}) is equivalent to minimizing $\beta$ subject to (\ref{eqn:icc}), (\ref{eqn:pc}). Rearranging gives that (\ref{eqn:icc}), (\ref{eqn:pc}) can be rewritten as
\begin{align}
    \beta &\geq \Delta + \alpha \label{eqn:icc2}\\
    \beta &\geq \Delta - U(\pi_0)/(U(\pi_1)-U(\pi_0)) \label{eqn:pc2}
\end{align}
Thus, the principal's problem is equivalent to minimizing $\beta$ subject to (\ref{eqn:icc2}), (\ref{eqn:pc2}). Call this optimization problem $\mathbf{P1}$.

%Our next step is to consider the problem of minimizing $\alpha$ subject to (\ref{eqn:pc2}). Call this optimization problem $\mathbf{P2}$. By the assumptions corresponding to (\ref{eqn:alph}), (\ref{eqn:beth}), we must have that any optimal solution to $\mathbf{P2}$ is also an optimal solution to $\mathbf{P1}$. But applying the generalized Neyman-Pearson Lemma (see for instance Theorem 3.6.1 of \cite{lehmann2022testing}) implies that (\ref{eqn:lrt}) is optimal for $\mathbf{P2}$ and hence also for $\mathbf{P1}$.
We next consider minimizing $\alpha$ subject to (\ref{eqn:pc2}). Call this optimization problem $\mathbf{P2}$. By the assumptions corresponding to (\ref{eqn:alph}), (\ref{eqn:beth}), the generalized Neyman-Pearson Lemma (see for instance Theorem 3.6.1 of \cite{lehmann2022testing}) implies that (\ref{eqn:lrt}) is optimal for $\mathbf{P2}$ when $\kappa \geq 0$ is chosen such that (\ref{eqn:beth}) holds. (Such a value for $\kappa$ exists because the relevant distributions are continuous.) This means (\ref{eqn:lrt}) is an optimal solution to $\mathbf{P2}$ where (\ref{eqn:alph}), (\ref{eqn:beth}) hold, which implies that this optimal solution to $\mathbf{P2}$ is also an optimal solution to $\mathbf{P1}$.
\end{proof}

The optimal decision region (\ref{eqn:lrt}) corresponds to a likelihood ratio test (LRT). And though the above result assumes $T,\pi_0,\pi_1$ are fixed, our next result shows that the optimal contract still retains this LRT structure even when we also optimize over the decision variables $T,\pi_0,\pi_1$.

\begin{theorem}
\label{thm:lrtopt}
If we have that $H_0 \neq H_1$, then an optimal solution $T^*, \pi_0^*, \pi_1^*, \mathcal{R}^*$ exists for the contract design (\ref{eqn:cost})--(\ref{eqn:pc}). The optimal payments are
\begin{align}
\pi_0^* &= \textstyle U^{-1}\Big(\gamma_a^{-T^*}\cdot\frac{-\alpha(\mathcal{R}^*)}{\beta(\mathcal{R}^*) - \alpha(\mathcal{R}^*)}\cdot(J_1^A-J_0^A)\Big)\label{eqn:pi0star}\\
\pi_1^* &= \textstyle U^{-1}\Big(\gamma_a^{-T^*}\cdot\frac{1-\alpha(\mathcal{R}^*)}{\beta(\mathcal{R}^*) - \alpha(\mathcal{R}^*)}\cdot(J_1^A-J_0^A)\Big)\label{eqn:pi1star}
\end{align}
Furthermore, the optimal $\mathcal{R}^*$ takes the form of (\ref{eqn:lrt}) with $T^*$.
\end{theorem}

\begin{proof} 
For the contract design problem, suppose we solve the contract design problem (\ref{eqn:cost})--(\ref{eqn:pc}) by optimizing over the decision variables in the following order:
\begin{equation}
\label{eqn:opt1}
\min_{T\in\mathbb{Z}_+}\,\min_{\mathcal{R}\subseteq\mathbb{R}^{T\cdot p}}\ \min_{\pi_0,\pi_1 : \pi_1 \geq \pi_0}
\end{equation}
The third minimization is interpreted as optimizing $\pi_0,\pi_1$ subject to a fixed $T$ and $\mathcal{R}$. A standard result (Proposition 4.4 of \cite{Laffont}) shows  the solution to this third optimization is 
\begin{figure*}[!t]
\normalsize
\begin{equation}
\label{eqn:uvw}
    U_i = \begin{bmatrix}
C & 0 & \dots & 0 \\
C(A+BK_i) & C & \dots & 0 \\
C(A+BK_i)^2 & C(A+BK_i) & \dots & 0 \\
\vdots & \vdots & \ddots & \vdots \\
C(A+BK_i)^{T-1} & C(A+BK_i)^{T-2} & \dots & C
\end{bmatrix}\qquad  V_i = \begin{bmatrix} C(A+BK_i)^{\hphantom{1}} \\ C(A+BK_i)^2 \\ C(A+BK_i)^3 \\ \vdots \\ C(A+BK_i)^T \end{bmatrix}
\end{equation}
\hrulefill % Optional: adds a horizontal line to separate the math from the text
\vspace*{4pt}
\end{figure*}
\begin{align}
\pi_0^*(\mathcal{R}) &= \textstyle U^{-1}\Big(\gamma_a^{-T}\cdot\frac{-\alpha(\mathcal{R})}{\beta(\mathcal{R}) - \alpha(\mathcal{R})}\cdot(J_1^A-J_0^A)\Big)\\
\pi_1^*(\mathcal{R}) &= \textstyle U^{-1}\Big(\gamma_a^{-T}\cdot\frac{1-\alpha(\mathcal{R})}{\beta(\mathcal{R}) - \alpha(\mathcal{R})}\cdot(J_1^A-J_0^A)\Big)
\end{align}
Optimality of (\ref{eqn:pi0star}), (\ref{eqn:pi1star}) follows by finishing problem (\ref{eqn:opt1}) by optimizing over $\mathcal{R}$ and then $T$. This optimal solution exists because there exists $\mathcal{R}$ such that $\beta(\mathcal{R}) > \alpha(\mathcal{R})$ since we have $H_0 \neq H_1$ (see for instance Corollary 3.2.1 of \cite{lehmann2022testing}).

Next, let $J(T,\pi_0,\pi_1,\mathcal{R})$ be the objective function (\ref{eqn:cost}) in the contract design problem. We note that
\begin{multline}
    \min_{\mathcal{R}\subseteq\mathbb{R}^{T\cdot p}} \big\{J(T^*,\pi_0^*,\pi_1^*,\mathcal{R}) : (\ref{eqn:icc}), (\ref{eqn:pc}) \big\} \geq \\
\min_{\substack{T\in\mathbb{Z}_+, \mathcal{R}\subseteq\mathbb{R}^{T\cdot p}\\\pi_0,\pi_1 : \pi_1 \geq \pi_0}} \big\{J(T,\pi_0,\pi_1,\mathcal{R}) : (\ref{eqn:icc}), (\ref{eqn:pc})\big\}.
\end{multline}
Since equality occurs when setting $\mathcal{R} = \mathcal{R}^*$ on the left-hand side, this means that solving the optimization problem
\begin{equation}
    \min_{\mathcal{R}\subseteq\mathbb{R}^{T\cdot p}} \big\{J(T^*,\pi_0^*,\pi_1^*,\mathcal{R}) : (\ref{eqn:icc}), (\ref{eqn:pc}) \big\}
\end{equation}
is equivalent to solving the original contract design problem (\ref{eqn:cost})--(\ref{eqn:pc}). This equivalent optimization problem matches the setup of Proposition \ref{prop:optimality}, and so the result follows if we can show the existence of an $\mathcal{R}$ satisfying (\ref{eqn:alph}), (\ref{eqn:beth}). Note that some algebra shows $- U(\pi_0^*)/(U(\pi_1^*)-U(\pi_0^*)) = \alpha(\mathcal{R}^*)$
%    U(\pi_1^*) - U(\pi_0^*) = \textstyle\gamma_a^{-T^*}\cdot\frac{1}{\beta(\mathcal{R}^*) - \alpha(\mathcal{R}^*)}\cdot(J_1^A-J_0^A).
and for (\ref{eqn:Delta}) that $\Delta = \beta(\mathcal{R}^*) - \alpha(\mathcal{R}^*)$. We finally argue this implies the existence of an $\mathcal{R}$ satisfying (\ref{eqn:alph}), (\ref{eqn:beth}): Indeed, choosing $\mathcal{R} = \mathcal{R}^*$ satisfies (\ref{eqn:alph}), (\ref{eqn:beth}).
\end{proof}

\subsection{Limited Liability Constraints}
In some applications, it can be preferred to impose an additional constraint of the form $\pi_0 \geq 0$ on the contract design problem (\ref{eqn:cost})--(\ref{eqn:pc}). Such a constraint is known as a \emph{limited liability} constraint, and it imposes a requirement that the agent can never receive a negative payment (i.e., the agent can never be fined). Note that without this constraint, Theorem \ref{thm:lrtopt} shows that in general $\pi_0$ will be negative, indicating that the agent would be fined for a state trajectory that looks like that which would be likely under low effort.


\begin{theorem}
\label{thm:lrtopt_ll}
If we have that $H_0 \neq H_1$, then an optimal solution $T^*, \pi_0^*, \pi_1^*, \mathcal{R}^*$ exists for the contract design problem (\ref{eqn:cost})--(\ref{eqn:pc}) with a limited liability constraint $\pi_0 \geq 0$. The optimal payments are given by
\begin{align}
\pi_0^* &= 0\\
\pi_1^* &= \textstyle U^{-1}\Big(\gamma_a^{-T^*}\cdot\frac{1}{\beta(\mathcal{R}^*) - \alpha(\mathcal{R}^*)}\cdot(J_1^A-J_0^A)\Big)\label{eqn:pi1star_ll}
\end{align}
Furthermore, the optimal $\mathcal{R}^*$ has the form of (\ref{eqn:lrt}) with $T^*$.
\end{theorem}

We omit the proof because it is roughly similar to the proof of Theorem \ref{thm:lrtopt}, with the main change being that Proposition \ref{prop:optimality} has to be changed to account for the fact that only the incentive compatibility constraint is binding under limited liability whereas both the incentive compatibility and participation constraints are binding in the original result.

\subsection{Decision to Induce Effort}
\label{sec:dieff}

After the principal has solved the contract design problem (\ref{eqn:cost})--(\ref{eqn:pc}) with or without limited liability constraints, which assumes the contract induces high effort from the agent, then the principal must decide whether it is beneficial from a cost perspective to induce high effort. If the principal declines to induce high effort, then they can simply offer a contract that consists of zero payment (i.e., they can set $\pi_0 = \pi_1 = 0$ with the choice of $T$ and $\mathcal{R}$ irrelevant) \cite{Laffont}. The principal would then incur an expected cost of $J^P_0$. 

The decision to induce effort now depends upon whether the following inequality is satisfied:
\begin{equation}
\label{eqn:eqndieff}
    J_1^P + \gamma_p^{T^*} \cdot \big(\pi_0^* \cdot (1-\beta(\mathcal{R}^*) + \pi_1^* \cdot \beta(\mathcal{R}^*)\big) \leq J_0^P,
\end{equation}
where $T^*,\pi_0^*,\pi_1^*,\mathcal{R}^*$ solve the contract design problem (\ref{eqn:cost})--(\ref{eqn:pc}) with or without liability constraints. The left-hand side of this inequality is the expected cost to the principal under an optimal contract that induces high effort, and the right-hand side is the expected cost to the principal when they offer no payment to the agent. If this inequality holds, then it is optimal to induce high effort and the principal offers $T^*,\pi_0^*,\pi_1^*,\mathcal{R}^*$ as the contract to the agent. If this inequality does not hold, then it is optimal to decline to induce high effort and the principal offers zero payment to the agent.

\section{COMPUTING THE OPTIMAL CONTRACT}
\label{sec:ctoc}
In this section, we provide an approach to design an optimal contract. We derive the likelihood ratio distribution, and  we provide an algorithm to compute the optimal contract. 


\subsection{Derivation of Likelihood Ratio Test}
Since Theorems \ref{thm:lrtopt} and \ref{thm:lrtopt_ll} show that the LRT provides an optimal decision region, our next step is to derive the LRT. By a slight abuse of notation, interpret $\mathbf{Y}_T \in \mathbb{R}^{T\cdot p}$ as a vector defined by stacking the observations $y_1,\ldots,y_T$. We first derive the distribution of $\mathbf{Y}_T$.

\begin{proposition}
\label{prop:ytdist}
    Under the state-feedback controller $K_i$ for $i = 0$ or $i = 1$, we have $\mathbf{Y}_T \sim \mathcal{N}(M_i, S_i)$ for $M_i = U_i\boldsymbol{\mu}_W + V_i\mu_0$ and $S_i = U_i\mathbf{\Sigma}_WU_i^\top + V_i\Sigma_0V_i^\top + \mathbf{\Sigma}_E$, with $U_i\in\mathbb{R}^{(T\cdot p)\times(T\cdot n)}$ and $V_i \in \mathbb{R}^{(T\cdot p)\times n}$ given in (\ref{eqn:uvw}), and where $\boldsymbol{\mu}_W \in \mathbb{R}^{T\cdot n}$ is defined by stacking $\mu_w$ a total of $T$ times, $\mathbf{\Sigma}_W = \textrm{blkdiag}(\Sigma_w, \ldots, \Sigma_w) \in \mathbb{R}^{(T\cdot n)\times(T\cdot n)}$ where $\textrm{blkdiag}(\cdot)$ specifies a block diagonal matrix, and $\mathbf{\Sigma}_E = \textrm{blkdiag}(\Sigma_e, \ldots,\Sigma_e) \in \mathbb{R}^{(T\cdot p)\times(T\cdot p)}$.
\end{proposition}

\begin{proof}
    Let $\mathbf{W}_T \in \mathbb{R}^{T\cdot n}$ be defined by stacking the random variables $w_0,\ldots,w_{T-1}$, and let $\mathbf{E}_T \in \mathbb{R}^{T\cdot p}$ be defined by stacking the random variables $\epsilon_0,\ldots,\epsilon_{T-1}$. Under the assumptions in our model, it is well known that $x_{t} = (A+BK_i)^tx_0 + \sum_{\tau = 0}^{t-1}(A+BK_i)^{t-\tau-1}w_\tau$, implying
    \begin{equation}
        \mathbf{Y}_T = U_i\mathbf{W}_T + V_ix_0 + \mathbf{E}_T.\label{eqn:lincom}
    \end{equation} 
    Since $x_0, \mathbf{W}_T, \mathbf{E}_T$ are jointly Gaussian and independent, they are all together jointly Gaussian. Hence, any linear combination such as (\ref{eqn:lincom}) is also jointly Gaussian. This means we completely know the distribution of $\mathbf{Y}_T$ if we calculate its mean and covariance. By linearity of expectation, we have $\mathbb{E}(\mathbf{Y}_T) = U_i\mathbb{E}(\mathbf{W}_T) + V_i\mathbb{E}(x_0) + \mathbb{E}(\mathbf{E}_T) = 
        U_i\boldsymbol{\mu}_W + V_i\mu_0$. If we denote the covariance using the notation $\textrm{Cov}(\cdot)$, then $\textrm{Cov}(\mathbf{Y}_T) = \textrm{Cov}(U_i\mathbf{W}_T) + \textrm{Cov}(V_ix_0) + \textrm{Cov}(\mathbf{E}_T) = 
        U_i\mathbf{\Sigma}_WU_i^\top + V_i\Sigma_0V_i^\top + \mathbf{\Sigma}_E$, where we used independence of $x_0,\mathbf{W}_T,\mathbf{E}_T$ in the first equality. The result follows.
\end{proof}

Because jointly Gaussian random variables belong to the exponential family, it is easier to implement the LRT using the log-likelihood ratio. Our next result characterizes the distribution of the log-likelihood ratio.

\begin{theorem}\label{thm:distribution}
Let $L_0(\mathbf{Y}_T)$ be the likelihood of $\mathbf{Y}_T$ under $H_0$, and let $L_1(\mathbf{Y}_T)$ be the likelihood of $\mathbf{Y}_T$ under $H_1$. The log-likelihood ratio $LLR = \log L_1(\mathbf{Y}_T) - \log L_0(\mathbf{Y}_T)$ follows a generalized chi-squared distribution, expressed as a linear combination of independent non-central chi-squared variables and a standard normal random variable:
\begin{align}
    LLR \mid H_0 &\sim \textstyle\sum_{j \in \mathcal{K}_0} w_{j} \chi^2_1(\nu_{j}^2) + \sigma_{0} Z_0 + C_0, \\
    LLR \mid H_1 &\sim \textstyle\sum_{j \in \mathcal{K}_1} \bar{w}_{j} \chi^2_1(\bar{\nu}_{j}^2) + \sigma_{1} Z_1 + C_1,
\end{align}
where $Z_0, Z_1 \sim \mathcal{N}(0,1)$, and $\mathcal{K}_0, \mathcal{K}_1$ are the index sets of dimensions with non-zero weights. The distribution parameters are determined by the mean shift $D = M_1 - M_0$ and the covariance matrices $S_0, S_1$. Define the core symmetric matrices with their respective eigenvalue decompositions: $\Phi_0 = S_0^{1/2} S_1^{-1} S_0^{1/2} = \Psi_0 \Lambda \Psi_0^\top$ and $\Phi_1 = S_1^{1/2} S_0^{-1} S_1^{1/2} = \Psi_1 \bar{\Lambda} \Psi_1^\top$. The matrices $\Psi_0, \Psi_1 \in \mathbb{R}^{(T\cdot p) \times (T\cdot p)}$ are orthogonal matrices whose columns consist of the eigenvectors of $\Phi_0$ and $\Phi_1$, respectively. $\Lambda, \bar{\Lambda} \in \mathbb{R}^{(T\cdot p) \times (T\cdot p)}$ are the diagonal matrices containing the corresponding eigenvalues.
The generalized chi-squared weights and rotated linear coefficients are derived directly from these eigenvalues and eigenvectors: $w_j = (1 - \lambda_j)/2$ and $\bar{w}_j = (\bar{\lambda}_j - 1)/2$ with $\lambda_j$ and $\bar{\lambda}_j$ being the diagonal elements of $\Lambda$ and $\bar{\Lambda}$. The other parameters are: $b_0 = \Psi_0^\top S_0^{1/2} S_1^{-1} D$, $b_1 = \Psi_1^\top S_1^{1/2} S_0^{-1} D$, $nu_{j} = \frac{b_{0,j}}{2w_{j}}$, $\bar{\nu}_{j} = \frac{b_{1,j}}{2\bar{w}_{j}}$, $\sigma_0^2 = \sum_{j \notin \mathcal{K}_0} b_{0,j}^2$, $\sigma_1^2 = \sum_{j \notin \mathcal{K}_1} b_{1,j}^2$, $C_0 = ( \log |S_{0}|/|S_{1}| - D^\top S_{1}^{-1} D )/2 - \sum_{j \in \mathcal{K}_0} w_{j} \nu_{j}^2$, and $C_1 =( \log |S_{0}|/|S_{1}| + D^\top S_{0}^{-1} D )/2 - \sum_{j \in \mathcal{K}_1} \bar{w}_{j} \bar{\nu}_{j}^2$.
\end{theorem}

\begin{proof}
The log-likelihood ratio simplifies to $LLR = \big(\log |S_{0}|/|S_{1}| + (\mathbf{Y}_T - M_0)^\top S_0^{-1}(\mathbf{Y}_T - M_0) + 
    - (\mathbf{Y}_T - M_1)^\top S_1^{-1}(\mathbf{Y}_T - M_1) \big)/2$. By Proposition \ref{prop:ytdist}, the stacked observation vector is distributed as $\mathbf{Y}_T \sim \mathcal{N}(M_i, S_i)$ under hypothesis $i \in \{0, 1\}$. We determine the distribution of the log-likelihood ratio by expanding it into a multivariate quadratic form and applying the exact standardization and diagonalization method detailed in Das and Geisler \cite{das2020methods}.

\vspace{1ex}
\noindent \textbf{Case 1: Under the Null Hypothesis ($H_0$).} \\
Here, $\mathbf{Y}_T \sim \mathcal{N}(M_0, S_0)$. Following \cite{das2020methods}, we first apply the whitening transformation $\mathbf{z} = S_0^{-1/2}(\mathbf{Y}_T - M_0) \sim \mathcal{N}(0, I)$. This transforms $S$ to a new quadratic form $\tilde{q}(\mathbf{z}) = \mathbf{z}^\top \tilde{Q}_0 \mathbf{z} + \tilde{q}_0^\top \mathbf{z} + c_{0}$, where the standardized coefficients are computed as $\tilde{Q}_0 = (I - S_0^{1/2} S_1^{-1} S_0^{1/2})/2 = (I - \Phi_0)/2$, $\tilde{q}_0 = S_0^{1/2} S_1^{-1} D$, and $c_{0} = ( \log |S_0|/|S_1| - D^\top S_1^{-1} D )/2$. Because $\Phi_0$ is orthogonally diagonalizable as $R_0 \Lambda R_0^\top$, we can substitute this to decompose $\tilde{Q}_0$ as $\tilde{Q}_0 = (I - \Psi_0 \Lambda \Psi_0^\top)/2 = \Psi_0 [ (I - \Lambda)/2 ] \Psi_0^\top \equiv \Psi_0 W \Psi_0^\top$, where $W$ is diagonal with  $w_j = (1 - \lambda_j)/2$. We define the rotated standard normal vector $\mathbf{y} = \Psi_0^\top \mathbf{z} \sim \mathcal{N}(0, I)$. Substituting $\mathbf{z} = \Psi_0 \mathbf{y}$ decouples the quadratic form into a sum of independent scalar variables:
\begin{align}
    \tilde{q}(\mathbf{y}) &= (\Psi_0 \mathbf{y})^\top (\Psi_0 W \Psi_0^\top) (\Psi_0 \mathbf{y}) + \tilde{q}_0^\top (\Psi_0 \mathbf{y}) + c_{0} \nonumber \\
    &= \mathbf{y}^\top W \mathbf{y} + b_0^\top \mathbf{y} + c_{0} \nonumber \\
    &= \textstyle\sum_{j=1}^{T\cdot p} \left( w_j y_j^2 + b_{0,j} y_j \right) + c_{0},
\end{align}
where $b_0 = \Psi_0^\top \tilde{q}_0$. We now partition this sum into dimensions with non-zero weights ($j \in \mathcal{K}_0$) and zero weights ($j \notin \mathcal{K}_0$). For dimensions where $w_j \neq 0$, we complete the square: $w_j y_j^2 + b_{0,j} y_j = w_j ( y_j + b_{0,j}/(2w_j) )^2 - b_{0,j}^2/(4w_j) \sim w_j \chi^2_1(\nu_j^2) - w_j \nu_j^2$, where $\nu_j = b_{0,j}/(2w_j)$ defines the non-centrality parameter. For dimensions where $w_j = 0$, the quadratic term vanishes. The remaining linear terms aggregate into a single standard normal distribution: $\sum_{j \notin \mathcal{K}_0} b_{0,j} y_j \sim \mathcal{N}(0, \sum_{j \notin \mathcal{K}_0} b_{0,j}^2) \equiv \sigma_0 Z_0$. Finally, substituting these independent random variables back into the sum and absorbing the constants yields the distribution for $H_0$: $LLR \mid H_0 \sim \sum_{j \in \mathcal{K}_0} w_j \chi^2_1(\nu_j^2) + \sigma_0 Z_0 + c_{0} - \sum_{j \in \mathcal{K}_0} w_j \nu_j^2$, where we get $C_0 = c_{0} - \sum_{j \in \mathcal{K}_0} w_j \nu_j^2$.

\vspace{1ex}
\noindent \textbf{Case 2: Under the Alternative Hypothesis ($H_1$).} \\
By symmetry, under $H_1$, $\mathbf{Y}_T \sim \mathcal{N}(M_1, S_1)$. The standardizing transformation $\bar{\mathbf{z}} = S_1^{-1/2}(\mathbf{Y}_T - M_1)$ maps $S$ into the quadratic form $\bar{q}(\bar{z}) = \bar{\mathbf{z}}^\top \tilde{Q}_1 \bar{\mathbf{z}} + \tilde{q}_1^\top \bar{\mathbf{z}} + c_{1}$ where the quadratic coefficients are defined as $\tilde{Q}_1 = (S_1^{1/2} S_0^{-1} S_1^{1/2} - I)/2 = (\Phi_1 - I)/2$, $\tilde{q}_1 = S_1^{1/2} S_0^{-1} D$, and $c_{1} = ( \log |S_0|/|S_1| + D^\top S_0^{-1} D )/2$. Applying the same direct decomposition and rotation decoupling as in the $H_0$ case algebraically produces the alternative weights $\bar{w}_j$, non-centralities $\bar{\nu}_j$, normal variance $\sigma_1^2$, and offset $C_1$, which concludes the proof.
\end{proof}

\begin{table}[t]
    \centering
    \caption{Mathematical Assumptions and Simulation Parameters}
    \label{tab:math_assumptions}
    \renewcommand{\arraystretch}{1.05} % Adds padding for mathematical equations
    \begin{tabular}{lcc}
        \hline
        \textbf{Parameter} & \textbf{Ex. 1: LFC} & \textbf{Ex. 2: Wellness} \\
        \hline
        $x_{k+1}$ & $Ax_k + Bu_k + \omega_k$ & $a x_k + b u_k + \omega_k$ \\
        $y_k$ & $Cx_k + \epsilon_k$ & $c x_k + \epsilon_k$ \\
        $x_k$ & $[x_g, x_t, x_p, x_i]^\top$ & $\text{scalar weight}$ \\
        $C$ & $I_{4 \times 4}$ & $1$ \\
        $x_0$ & $\mathcal{N}([0,0,0.1,0]^\top, 10^{-9}\cdot I)$ & $\mathcal{N}(65, 0)$\\% (deterministic) \\
        $w_k$ & $\mathcal{N}(0, 0.01\cdot I)$ & $\mathcal{N}(0.817, 0.5)$ \\
        $\epsilon_k$ & $\mathcal{N}(0, 0.01 \cdot I)$ & $\mathcal{N}(0, 0.01)$ \\
        $u_0$ & $K_0 x_k (\lambda = 0.2)$ & $\delta_0x_k(\delta_0 = 70)$ \\
        $u_1$ & $K_1 x_k (\lambda = 0.05)$ & $\delta_1x_k(\delta_1 = 110)$  \\
        $J^A_1 - J^A_0$ & 0.1 & 5\\
        Time Step & $0.1s$ & 1 week \\
        \hline
    \end{tabular}
\end{table}
\subsection{Computational Algorithm}

Our characterization of the optimal contract structure in Section \ref{sec:soc} is significant because  it allows us to convert solving the contract design problem (\ref{eqn:cost})--(\ref{eqn:pc}), which is a complex, infinite dimensional-optimization problem, into two line searches, one of which is nested within the other.

For the contract design problem (\ref{eqn:cost})--(\ref{eqn:pc}), define the decision region $R(T,\eta) = \{LLR \geq \eta\}$ where $LLR$ is as defined in Theorem \ref{thm:distribution}, and define the payments
\begin{align}
\pi_0(T,\eta) &= \textstyle U^{-1}\Big(\gamma_a^{-T}\cdot\frac{-\alpha(\mathcal{R}(T,\eta))}{\beta(\mathcal{R(T,\eta)}) - \alpha(\mathcal{R}(T,\eta))}\cdot(J_1^A-J_0^A)\Big)\\
\pi_1(T,\eta) &= \textstyle U^{-1}\Big(\gamma_a^{-T}\cdot\frac{1-\alpha(\mathcal{R}(T,\eta))}{\beta(\mathcal{R(T,\eta)}) - \alpha(\mathcal{R(T,\eta)})}\cdot(J_1^A-J_0^A)\Big)
\end{align}
Then we can solve the contract design problem (\ref{eqn:cost})--(\ref{eqn:pc}) using the following algorithm:
\begin{enumerate}
    \item For $t \in \{1,\ldots,\overline{T}\}$,
    \begin{enumerate}
        \item For $\eta \in \mathbb{R}$, set $J(T,\eta) := J_1^P + \gamma_p^{T} \cdot (\pi_0(T,\eta) \cdot (1-\beta(\mathcal{R}(T,\eta)) + \pi_1(T,\eta) \cdot \beta(\mathcal{R}(T,\eta))$;
        \item Set $\eta^*(T) := \arg\min_{\eta \in \mathbb{R}}J(T,\eta)$;
    \end{enumerate}
    \item Set $T^* := \arg\min_{t\in\{1,\ldots,\overline{T}\}} J(T,\eta^*(T))$;
    \item Set $\eta^* := \eta^*(T^*)$, set $\pi_0^* := \pi_0(T^*,\eta^*)$, and set $\pi_1^* := \pi_1(T^*,\eta^*)$;
\end{enumerate}
Though steps 1.a and 1.b are infinite-dimensional, numerically we can consider a bounded range for $\eta$ and perform the calculations at a finite number of points within this bounded range. Also, note that $\alpha,\beta$ can be computed using Imhof's method, via Davies's algorithm \cite{Imhof1961, Davies1980, Das2025}. In our experiments detailed in the next section, we use the Python package \texttt{chi2comb} \cite{horta2021chi2comb} for the computation.


The contract design problem (\ref{eqn:cost})--(\ref{eqn:pc}) with a limited liability constraint $\pi_0 \geq 0$ can be solved with the algorithm given above but with the change that we use
\begin{align}
\pi_0(T,\eta) &= 0\\
\pi_1(T,\eta) &= \textstyle U^{-1}\Big(\gamma_a^{-T}\cdot\frac{1}{\beta(\mathcal{R(T,\eta)}) - \alpha(\mathcal{R(T,\eta)})}\cdot(J_1^A-J_0^A)\Big)
\end{align}
as our payment amounts in the algorithm.

\section{APPLICATION EXAMPLES}
\label{sec:ae}

\begin{figure}[t]
    \centering
    \includegraphics[width=0.45\textwidth]{fig/0401LFC.png}
    \caption{Optimal Contract Parameters for LFC Example.}
    \label{fig:optimal_lfc}
\end{figure}

We apply our framework to  applications from power systems and health care. We focus on the contract design problem (\ref{eqn:cost})--(\ref{eqn:pc}) with limited liability constraints (i.e., the agent will not accept being fined). In both examples, we use $U(\pi) = \sqrt{\pi}$ as the agent's utility function, which satisfies $U' > 0$ and $U'' < 0$ for nonnegative payments. We interpret $\pi$ as having units of US Dollars. We do not consider whether to induce effort as discussed in Section \ref{sec:dieff}, because this decision uses a simple inequality (\ref{eqn:eqndieff}) and because the answer is highly sensitive to the choice for the value of $J_1^P-J_0^P$.

In both examples, we also assume that the agent's discount factor is strictly lower than the principal's (i.e., $\gamma_a < \gamma_p$). This is expected if the agent has a higher return on investments than the principal. It is reasonable in our examples where the principal is a large entity, like an Independent System Operator (ISO) or a health insurance company, which would invest more conservatively than a smaller entity. In our simulations, we set $\gamma_a = 0.995$ and $\gamma_p = 0.998$. 
 
\subsection{Load Frequency Control (LFC) Model}
Load frequency control (LFC) \cite{Tan2010} is important for stabilizing the frequency in a power system. Here, the principal (e.g., the ISO) offers a performance-based contract to incentivize the agents (e.g., power generators) to adopt a high-performance control policy ($K_1$) for frequency regulation. When the ISO cannot directly inspect the generator's internal control settings, a moral hazard problem arises. The contract design we suggest enables the principal to infer the internal settings of the generator from a partial, noisy history of the system states, thereby incentivizing the implementation of high-performance control.

% We discretized the single-area LFC system with a sampling period of $T = 0.1s$ using the Zero-Order Hold (ZOH) method\cite{Tan2010}. The system is represented by the following discrete-time state-space equation:
% \begin{equation}
% x_{k+1} = Ax_k + Bu_k + \omega_k,\quad y_k = Cx_k + \epsilon_k
% \end{equation}
% where $x_k = [x_g, x_t, x_p, x_i]^\top$ is the state vector consisting of governor power, turbine power, frequency deviation, and the integral of frequency deviation. We also assumed $J^A_1 - J^A_0 = 0.1,  w_k \sim N(\mathbf{0}, 10^{-9}\cdot I), C = I$ and $\epsilon_k \sim N(0, 0.01\cdot I)$. Based on the benchmark parameters $T_G = 0.08, T_T = 0.3, T_P = 20$, $K_P = 120$, and $R =2.4$, the discrete-time system matrices are computed as follows:
We discretize a single-area LFC system with a sampling period of $T = 0.1s$ using Zero-Order Hold (ZOH) \cite{Tan2010}. The state-space distributions and system variables are summarized in Table \ref{tab:math_assumptions}. Based on benchmark parameters $T_G = 0.08, T_T = 0.3, T_P = 20, K_P = 120$, and $R =2.4$, the discrete-time system matrices are
\begin{equation*}
A = \begin{bmatrix}
0.287 & 0 & 0 & 0 \\
0.156 & 0.717 & 0 & 0 \\
0.061 & 0.509 & 0.995 & 0 \\
0.002 & 0.027 & 0.100 & 1
\end{bmatrix}\qquad 
B = \begin{bmatrix}
0.713 \\ 0.127 \\ 0.029 \\ 0.001
\end{bmatrix}.
\end{equation*}
We simplify the IMC-PID tuning from \cite{Tan2010} for the agent's controller: By omitting the disturbance rejection filter $Q_d(s)$, we choose $K_0$ (low effort) and $K_1$ (high effort) by setting the closed-loop time constant $\lambda$ to $0.2$ and $0.05$, respectively. This results in $K_0 = [0, 0.0002, -0.0856, -0.0139]$ and $K_1 = [0, -2.2033, -0.6932, -0.0556]$. 

Figure \ref{fig:optimal_lfc} shows the optimal transfer $\pi_1^*$, decision threshold $\ln \eta^*$, and expected cost as a function of the contract length $T$. We assume the agent maintains a chosen control setting over a physical operating cycle of 30 seconds, with the possibility of switching their controller every 30 seconds. Accordingly, we use a maximum observation horizon of 30 seconds ($\overline{T} = 300$). The globally optimal observation length is attained at $T^* = 18$ ($1.8$ seconds), where the principal’s expected cost is minimized at $30.0923$, requiring an optimal incentive payment of $\pi_1^* = \$31.6495$.

%Assuming a maximum practical history length of $30$ seconds ($T = 300$), the global optimal contract length is achieved at $T^* = 18$ ($1.8$ seconds). At this evaluation horizon, the principal's expected cost is minimized at $0.0120$, requiring an optimal incentive payment of $\pi^*_1 = 0.0127$.
% shows the changes in optimal transfer $\pi_1$, optimal threshold($\ln \eta$), and the minimum cost of the principal by change of the contract length along with the global optimal contract. Assume the physical maximum length of the history is $30s(T = 300)$. Then the global optimal $T_0$ is at $T^*_0 = 18 (1.8s)$ with cost $0.0120$ and $\pi^*_1 = 0.0127$.

\subsection{Wellness Incentive Program}
Imagine a wellness program where a health insurance company provides payments to encourage weight loss over 200 weeks. The Mifflin-St Jeor equation \cite{mifflin1990new,aswani2019behavioral} gives the coefficients of an affine LTI system to describe weight $x_t$ (kg) dynamics under modest amounts of weight change, where the coefficients depends upon an individual's age, starting weight, height, and gender. For a 30-year-old female ($h = 160.0$ cm, $x_0 = 65.0$ kg), the dynamics with a time step of one-week are $    x_{t+1} = 0.999 x_t -6.429 u_t + 1.286\times 10^{-4} f_t-0.618 + z_t$,  where $z_t$ is a random (biologically-driven) fluctuation, $f_t$ is caloric intake (we set $f_t \equiv 1600$ kcal), and $u_t$ is physical activity in units of steps. We modeled the agent's effort as a state-dependent feedback control $u_t = \delta_i x_t$, where $\delta_1 = 110$ represents a high-effort regimen and $\delta_0 = 70$ represents baseline activity. Figure \ref{fig:optimal_wellness} illustrates the trade-offs involved in designing the wellness incentive program. The global minimum of the insurance company's cost occurs with a contract length of 80 weeks, where the cost is 121.58 and optimal $\pi_1$ is \$217.02.

\begin{figure}[t]
    \centering
    \includegraphics[width=0.45\textwidth]{fig/FinalWellness.png}
    \caption{Optimal Contract Parameters for Wellness Example.}
    \label{fig:optimal_wellness}
\end{figure}

% \begin{table}[htbp]
%     \centering
%     \caption{Optimal Parameters over Various Contract Lengths}
%     \label{tab:contract_parameters}
%     \begin{tabular}{ccccc}
%         \hline
%         \textbf{Week} & $\pi_1^*$ & \textbf{Cost} & $\ln \eta^*$ & $\beta^*-\alpha^*$ \\
%         \hline
%         1   &8683.27 &2579.51 &0.1 &0.0539 \\
%         4   &2213.42 &724.51 &0.2 & 0.1084  \\
%         8   &963.52 & 403.28& 0.2 & 0.1676\\
%         26  &349.77 & 176.13&0.3 &0.3044\\
%         52  & 237.34& 129.72& 0.4& 0.4206  \\
%         80* & 217.02&121.58&0.5& 0.5058\\
%         104 & 223.30& 127.46& 0.5& 0.5621 \\
%         \hline
%     \end{tabular}
% \end{table}
% \begin{figure}[htbp]
%     \centering
%     \includegraphics[width=0.45\textwidth]{fig/insurance_delta5.png}
%     \caption{Optimal Contract Length and Optimal Payment and Costs.}
%     \label{fig:optimal_threshold}
% \end{figure}



% For short contract lengths, the principal incurs significantly higher costs. This occurs because there is less time to deviate the weight trajectory from that of low effort, in other words, the test is less informative (lower $\beta - \alpha$) requiring a much larger payment $\pi_1$ to satisfy the participation constraint against stochastic metabolic noise. However, the difference in time preference of agent and principal ($\frac{\gamma_p}{\gamma_a^2}^{T_0}$) cancels out the effect of increasing $\beta-\alpha$, which yields optimal contract length at $80$.
\subsection{Results and Discussion}
Despite the different physical contexts of the two examples, our simulations reveal a functionally identical trade-off in the contract design. Figures \ref{fig:optimal_lfc} and \ref{fig:optimal_wellness} illustrate how the principal's expected cost and the optimal high-effort payment ($\pi_1^*$) evolve across varying contract evaluation horizons ($T$). In both scenarios, short evaluation horizons force the principal to incur higher expected costs. Over a short horizon, system stochasticity dominates the observed state trajectories, forcing the principal to offer a substantially larger payment $\pi_1^*$ to overcome uncertainty and make the contract feasible. As $T$ increases, the structural divergence between the high-effort and low-effort policies begins to overpower the noise ($w_k$, $\epsilon_k$). This effectively reduces the informational friction between the principal and the agent, reducing the required statistical risk premium. After a certain point, this statistical benefit is counterbalanced by the cost of delayed compensation driven by the discount factor asymmetry ($\gamma_a < \gamma_p$) as the exponential multiplier $(\gamma_p / \gamma_a^2)^T$ grows, creating the  U-shaped cost curves observed.

\section{CONCLUSION}
%In this paper, we introduced an optimal contract design framework for principal-agent problems characterized by pure moral hazard in stochastic, discrete-time LTI systems. We established a hypothesis testing-based incentive mechanism that enables the principal to infer unobservable agent actions from noisy, continuous state trajectories. We proved the likelihood ratio test-based contract is optimal contract by relating the bilevel optimization of the principal-agent problem to the Neyman-Pearson lemma, and we derived the exact distribution of the log-likelihood ratio. Furthermore, we identified the optimal payments for both high-effort and low-effort outcomes, accounting for scenarios both with and without limited liability constraints. Through our analysis and simulations, we demonstrated that the optimal contract evaluation horizon is governed by a trade-off of extending the observation window which allows a statistically more powerful test, and the cost of delayed compensations.

%In this paper, we introduced a model for incentive design under moral hazard in stochastic LTI systems. When the agent chooses between two linear state-feedback controllers, we showed that the optimal contract with two payment levels uses a likelihood ratio test to determine payment amounts. We finished by conducting numerical experiments with two examples, which demonstrated that the optimal contract evaluation horizon is governed by a trade-off of extending the observation window, which allows a statistically more powerful test, and the cost of delayed compensations. %Future work involves generalizing our results to nonlinear dynamics and controllers and non-quadratic costs, though we note that some such generalizations follow almost trivially. 

We studied incentive design for stochastic LTI systems with moral hazard. When the agent chooses between two controllers, we showed that the optimal contract with two payment levels uses a likelihood ratio test to determine payments. We finished by conducting numerical experiments with two examples, which demonstrated that the optimal contract evaluation horizon is governed by a trade-off of extending the observation window, which allows a statistically more powerful test, and the cost of delayed compensations. Future work involves generalizing our results to nonlinear dynamics and controllers and non-quadratic costs.%, though we note that some such generalizations follow almost trivially. 

\bibliographystyle{IEEEtran}
\bibliography{main}

\end{document}
